\documentclass[11pt,a4paper]{article}
\usepackage[utf8]{inputenc}
\usepackage{amsmath,amssymb,amsthm}
\usepackage{geometry}
\geometry{margin=1in}

\newtheorem{definition}{Definition}
\newtheorem{proposition}{Proposition}
\newtheorem{conjecture}{Conjecture}
\newtheorem{remark}{Remark}

\title{Holographic Entropy Duality: A Definition}
\author{Valhalla Research Team}
\date{January 2026}

\begin{document}
\maketitle

\begin{abstract}
We define an entropy duality transformation and explore its algebraic properties. The transformation has a fixed point at $S^* = 4\pi$, which corresponds to the Planck-mass black hole entropy. We emphasize that this is a \textbf{definition}, not a derivation from fundamental physics. The physical significance, if any, remains speculative.
\end{abstract}

\section{Definition}

\begin{definition}[Entropy Dual]
For entropy $S > 0$, define the entropy dual:
\begin{equation}
S' := \frac{(4\pi)^2}{S}
\end{equation}
\end{definition}

\section{Algebraic Properties}

\begin{proposition}[Involution]
The transformation is an involution:
\begin{equation}
(S')' = \frac{(4\pi)^2}{S'} = \frac{(4\pi)^2}{(4\pi)^2/S} = S
\end{equation}
\end{proposition}

\begin{proposition}[Fixed Point]
The unique fixed point is:
\begin{equation}
S^* = 4\pi
\end{equation}
Proof: $S^* = (4\pi)^2/S^* \Rightarrow (S^*)^2 = (4\pi)^2 \Rightarrow S^* = 4\pi$.
\end{proposition}

\begin{proposition}[Product Invariant]
For any entropy $S$:
\begin{equation}
S \cdot S' = (4\pi)^2
\end{equation}
\end{proposition}

\section{Physical Context}

\subsection{Planck-Mass Black Hole}

A black hole with mass equal to the Planck mass has Schwarzschild radius $r_s = 2l_P$ and entropy:
\begin{equation}
S_{\text{Planck BH}} = \frac{A}{4l_P^2} = \frac{4\pi r_s^2}{4l_P^2} = \frac{4\pi (2l_P)^2}{4l_P^2} = 4\pi
\end{equation}

This matches the fixed point $S^* = 4\pi$.

\subsection{Interpretation}

\begin{remark}
The fixed point at Planck-mass black hole entropy is suggestive of a UV-IR duality:
\begin{itemize}
    \item $S < 4\pi$: Sub-Planckian (UV) regime
    \item $S > 4\pi$: Super-Planckian (IR) regime
    \item $S = 4\pi$: Self-dual (Planck scale)
\end{itemize}
\end{remark}

\begin{conjecture}
The duality $S \cdot S' = (4\pi)^2$ may encode a UV-IR correspondence in quantum gravity, analogous to string T-duality where $R \cdot R' = \alpha'$.
\end{conjecture}

\section{Caution}

\begin{remark}[Definitional Nature]
This paper defines a mathematical transformation on entropy values. The physical significance is \textbf{speculative}. We do not claim:
\begin{itemize}
    \item This duality is derived from any fundamental theory
    \item There exists a physical mechanism implementing this transformation
    \item The analogy to T-duality is exact
\end{itemize}
\end{remark}

\section{Summary}

We have defined an entropy duality with algebraic properties:
\begin{center}
\begin{tabular}{|l|l|}
\hline
Property & Formula \\
\hline
Definition & $S' = (4\pi)^2/S$ \\
Involution & $(S')' = S$ \\
Fixed point & $S^* = 4\pi$ \\
Invariant & $S \cdot S' = (4\pi)^2$ \\
\hline
\end{tabular}
\end{center}

The fixed point coincides with Planck-mass black hole entropy, which is suggestive but not conclusive.

\begin{thebibliography}{9}
\bibitem{bekenstein1973} J.D. Bekenstein, \textit{Black holes and entropy}, Phys. Rev. D 7 (1973), 2333.
\bibitem{hawking1975} S.W. Hawking, \textit{Particle creation by black holes}, Comm. Math. Phys. 43 (1975), 199.
\bibitem{maldacena1998} J. Maldacena, \textit{The large N limit of superconformal field theories and supergravity}, Adv. Theor. Math. Phys. 2 (1998), 231.
\end{thebibliography}

\end{document}
